\documentclass[10pt,a4paper]{amsart}
\usepackage[margin=19mm]{geometry}
\usepackage{amsmath,amssymb,mathtools}
\usepackage[scr=rsfs,cal=euler]{mathalfa}
\usepackage{xfrac}
\usepackage[unicode,hidelinks,bookmarks=false]{hyperref}
\newcommand{\ie}{i.e.\ }
\newcommand{\cf}{cf.\ }
\newcommand{\resp}{respectively\ }
\newcommand{\Subsection}{\subsection}
\providecommand{\nobreakdash}{\nobreak\hbox{-}}
\setlength{\emergencystretch}{2em}
\multlinegap0pt
\usepackage{siunitx}
\usepackage{siunitx}
\usepackage{siunitx}
\usepackage{siunitx}
\usepackage{amssymb}
\usepackage{bm}
\usepackage{tikz}
\usepackage[shortlabels]{enumitem}
\usepackage{siunitx}
\usepackage{xcolor}
\newtheorem{lemma}{Lemma}
\newcommand{\Ael}{\mathcal{A}} % Caligraphic A
\newcommand{\Id}{\mathbb{I}} % Identity tensor (using \mathbb for blackboard bold)
\newcommand{\Thprod}[2]{( #1, #2 )_{\mathcal{T}_h}}
\newcommand{\divergence}{\mathrm{div}\,}
\newcommand{\grad}{\mathrm{grad}\,}
\newcommand{\pThprod}[2]{\langle #1, #2 \rangle_{\partial \mathcal{T}_{h}}}
\newcommand{\vect}[1]{\mathbf{#1}}
\title{Hybridizable discontinuous Galerkin methods for poroelastic wave propagation with symmetric stress approximation}
\author{Jeonghun J. Lee, Manuel A. Sanchez}
\date{Source: arXiv:2605.05451v1 (CC BY 4.0)}
\begin{document}
\maketitle

\noindent\emph{Source and licence.} Hybridizable discontinuous Galerkin methods for poroelastic wave propagation with symmetric stress approximation. Version arXiv:2605.05451v1 (CC BY 4.0), distributed by arXiv under the Creative Commons Attribution 4.0 International licence, http://creativecommons.org/licenses/by/4.0/, as recorded in the arXiv licence metadata for this exact version and shown on the abstract page https://arxiv.org/abs/2605.05451v1. Full licence notice: \texttt{licenses/arxiv-2605.05451-CC-BY-4.0.txt}.\par\smallskip

\noindent\textit{Editorial scope.} Counted unit: the article's lemma lemma (source0.tex lines 1309--1311). The article's complete original proof of it is delivered with it (source0.tex lines 1313--1367, one \texttt{proof} environment of 55 lines). No mathematics is written, repaired or re-derived: the statement and the proof are the article's own lines, in the article's own notation, and no retained line is substituted or re-flowed.  Retained uncounted as the source-local context the proof needs: lines 292--295; lines 306--311; lines 1261--1367.  Nothing was omitted from any retained span, so the package carries no omission record.  No retained line contains a citation, so the wrapper carries no bibliography and no bibliography entry.  No retained line contains a figure environment, an image-inclusion command or drawing code, so no image is delivered and the package declares no asset dependency.  Editorial formatting only, in addition: an AMS \texttt{amsart} preamble that restates the article's own declarations and macros used by the retained lines, namely the environments \texttt{lemma} on separate counters, together with the standard packages those lines need.  The retained environments print in this document as Lemma 1 in order of appearance, whereas the article prints the same items with the numbering of its own full document.  The complete native source of the article as distributed in the arXiv e-print for this exact version is bundled unchanged as \texttt{sources/199955/source0.tex}.
\begin{align}
	\label{eq:rho-coercivity}
	\rho_{11}\rho_{22}-\rho_{12}^2 \ge \rho_0 >0.
\end{align}

\begin{multline}
	\label{eq:A-positive}
	\int_{\Omega} \Ael r : r \,dx \ge 0 \text{ for } r \in L^2(\Omega; \mathbb{S}) 
    \\
    \text{ and } \quad r = 0 \in L^2(\Omega; \mathbb{S}) \quad \text{ if } \int_{\Omega} \Ael r : r \,dx = 0 
\end{multline}

Given $(\sigma_h^i, \vect{v}_{s,h}^i, \widehat{\vect{v}}_{s,h}^i, \vect{v}_{f,h}^i, p_h^i, \widehat{p}_h^i)$ and $\vect{f}^i$, $\vect{f}^{i+1}$, $g^i$, $g^{i+1}$, the Crank--Nicolson scheme seeks 
%
\begin{align*}
	(\sigma_h^{i+1}, \vect{v}_{s,h}^{i+1}, \widehat{\vect{v}}_{s,h}^{i+1}, \vect{v}_{f,h}^{i+1}, p_h^{i+1}, \widehat{p}_h^{i+1}) \in \Sigma_{h} \times V_{s,h} \times V_{f,h} \times Q_{h} \times \widehat{V}_{s,h}\times \widehat{Q}_{h}
\end{align*}
%
satisfying the following equations:  
%
\begin{subequations}
	\label{eq:fully-discrete-cn-eqs}
	\begin{align}
		\label{eq:fully-discrete-cn-eq1}
		\Thprod{\Ael(d_t \sigma_{h}^{i+1} + \alpha d_t{p}_{h}^{i+1} \Id)}{r} +\Thprod{\divergence r}{\vect{v}_{s,h}^{i+\frac 12}} -\pThprod{\widehat{\vect{v}}_{s,h}^{i+\frac 12}}{r \vect{n}} &= 0 ,
		\\ 
		\label{eq:fully-discrete-cn-eq2}		\Thprod{\rho_{11}d_t{\vect{v}}_{s,h}^{i+1}+\rho_{12}d_t{\vect{v}}_{f,h}^{i+1}}{\vect{w}_{s}} + \Thprod{\sigma_{h}^{i+\frac 12}}{\grad \vect{w}_{s}} &
        \\
        \notag 
        -\pThprod{\widehat{\sigma_{h} \vect{n}}^{i+\frac 12} }{\vect{w}_{s}} 
		&= {\Thprod{\vect{f}^{i+1/2}}{\vect{w}_{s}}}, 
		\\
		\label{eq:fully-discrete-cn-eq3}
		\Thprod{\rho_{12}d_t{\vect{v}}_{s,h}^{i+1}+\rho_{22}d_t{\vect{v}}_{f,h}^{i+1}}{\vect{w}_{f}} + \Thprod{\frac{\eta}{\kappa} \vect{v}_{f,h}^{i+\frac 12}}{\vect{w}_{f}} 
		&
        \\
        \notag 
        - \Thprod{p_{h}^{i+\frac 12} }{\divergence\vect{w}_{f}} +\pThprod{\widehat{p}_{h}^{i+\frac 12}}{\vect{w}_{f}\cdot\vect{n}} &= 0, 
		\\
		\label{eq:fully-discrete-cn-eq4}
		\Thprod{s_{0}d_t{p}_{h}^{i+1}}{q} + \Thprod{\Ael(d_t{\sigma}_{h}^{i+1}+\alpha d_t{p}_{h}^{i+1}\Id)}{\alpha q\Id} &
        \\
        \notag 
        - \Thprod{\vect{v}_{f,h}^{i+\frac 12}}{\grad q} 
		+\pThprod{\widehat{\vect{v}_{f,h}\cdot \vect{n}}^{i+\frac 12} }{q}
		&={\Thprod{g^{i+1/2} }{q}}, 
		\\
		\label{eq:fully-discrete-cn-eq5}
		\pThprod{\widehat{\sigma_h\vect{n}}^{i+\frac 12}  }{\widehat{\vect{w}}_{s}}
		=
		\pThprod{\sigma_h^{i+\frac 12} \vect{n} - \tau_{s} ({P_{\widehat{V}_s}} \vect{v}_{s,h}^{i+\frac 12} - \widehat{\vect{v}}_{s,h}^{i+\frac 12})}{\widehat{\vect{w}}_{s}} &= 0, 
		\\
		\label{eq:fully-discrete-cn-eq6}
		\pThprod{\widehat{\vect{v}_{f,h} \cdot \vect{n}}^{i+\frac 12} }{\widehat{q}}
		=
		\pThprod{\vect{v}_{f,h}^{i+\frac 12}\cdot \vect{n} + \tau_f (p_h^{i+\frac 12} - \widehat{p}_h^{i+\frac 12})}{\widehat{q}} &= 0
	\end{align}
\end{subequations}
for any $(r, \vect{w}_s, \vect{w}_f, q, \widehat{\vect{w}}_s, \widehat{q}) \in \Sigma_{h} \times V_{s,h} \times V_{f,h} \times Q_{h} \times \widehat{V}_{s,h}\times \widehat{Q}_{h}$. In the following lemma, we prove the well-posedness of the fully discrete scheme.
%
\begin{lemma}
	The system \eqref{eq:fully-discrete-cn-eqs} is well-posed.
\end{lemma}
%
\begin{proof}
	We first note that \eqref{eq:fully-discrete-cn-eqs} is a square linear system; thus, the uniqueness of solutions implies the existence of solutions. 
	Suppose then that $(\sigma_h^i, \vect{v}_{s,h}^i, \widehat{\vect{v}}_{s,h}^i, \vect{v}_{f,h}^i, p_h^i, \widehat{p}_h^i)$ and $g^i$, $g^{i+1}$, $\vect{f}^{i}$, $\vect{f}^{i+1}$ vanish. Then, \eqref{eq:fully-discrete-cn-eqs} reduces to 
	%
	% \begin{subequations}\label{eq:reduced-fully-discrete-cn-eqs}
		\begin{align*}
			% \label{eq:reduced-fully-discrete-cn-eq1}
			\Thprod{\Ael(\sigma_{h}^{i+1} + \alpha {p}_{h}^{i+1} \Id)}{r} +\frac{\Delta t}{2} (\Thprod{\divergence r}{\vect{v}_{s,h}^{i+1}} -\pThprod{\widehat{\vect{v}}_{s,h}^{i+1}}{r \vect{n}})& = 0 ,
			\\ 
			% \label{eq:reduced-fully-discrete-cn-eq2}
			\Thprod{\rho_{11}{\vect{v}}_{s,h}^{i+1} +\rho_{12}{\vect{v}}_{f,h}^{i+1}}{\vect{w}_{s}} + \frac{\Delta t}{2}\left(\Thprod{\sigma_{h}^{i+1}}{\grad \vect{w}_{s}} 
			- \pThprod{\widehat{\sigma_{h} \vect{n}}^{i+1} }{\vect{w}_{s}}\right)& = 0, 
			\\
			% \label{eq:reduced-fully-discrete-cn-eq3}
			\Thprod{\rho_{12}{\vect{v}}_{s,h}^{i+1}+\rho_{22}{\vect{v}}_{f,h}^{i+1}}{\vect{w}_{f}} + \frac{\Delta t}{2}\left(\Thprod{\frac{\eta}{\kappa} \vect{v}_{f,h}^{i+1}}{\vect{w}_{f}} 
			- \Thprod{p_{h}^{i+1} }{\divergence\vect{w}_{f}} +\pThprod{\widehat{p}_{h}^{i+1}}{\vect{w}_{f}\cdot\vect{n}}\right) &= 0, 
			\\
			% \label{eq:reduced-fully-discrete-cn-eq4}
			\Thprod{s_{0}{p}_{h}^{i+1}}{q} + \Thprod{\Ael({\sigma}_{h}^{i+1}+\alpha {p}_{h}^{i+1}\Id)}{\alpha q\Id} + \frac{\Delta t}{2} \left(-\Thprod{\vect{v}_{f,h}^{i+1}}{\grad q} 
			+ \pThprod{\widehat{\vect{v}_{f,h}\cdot \vect{n}}^{i+1}}{q}\right) &= 0,
			\\
			% \label{eq:reduced-fully-discrete-cn-eq5}
			\frac{\Delta t}{2}\pThprod{\widehat{\sigma_h \vect{n}}^{i+1} }{\widehat{\vect{w}}_{s}}=
			\frac{\Delta t}{2}\pThprod{\sigma_h^{i+1} \vect{n} - \tau_{s} (P_{\widehat{V}_s} \vect{v}_{s,h}^{i+1} - \widehat{\vect{v}}_{s,h}^{i+1})}{\widehat{\vect{w}}_{s}} &= 0, 
			\\
			% \label{eq:reduced-fully-discrete-cn-eq6}
			\frac{\Delta t}{2}\pThprod{\widehat{\vect{v}_{f,h}\cdot \vect{n}}^{i+1})}{\widehat{q}}=
			\frac{\Delta t}{2}\pThprod{\vect{v}_{f,h}^{i+1}\cdot \vect{n} + \tau_f (p_h^{i+1} - \widehat{p}_h^{i+1})}{\widehat{q}} &= 0 .
		\end{align*}
		% \end{subequations}
	%
	By taking the test functions $r = \sigma_h^{i+1}$, $\vect{w}_{s} = \vect{v}_{s,h}^{i+1}$, $\vect{w}_{f} = \vect{v}_{f,h}^{i+1}$, $q = p_h^{i+1}$, $\widehat{\vect{w}}_{s} = \widehat{\vect{v}}_{s,h}^{i+1}$, $\widehat{q}=\widehat{p}_{h}^{i+1}$ in the system above and adding all the equations, we obtain
	%
	\begin{align*}
		\Thprod{\Ael({\sigma}_{h}^{i+1} + \alpha {p}_{h}^{i+1} \Id)}{{\sigma}_{h}^{i+1} + \alpha {p}_{h}^{i+1} \Id)}
        \\
		+ \Thprod{\rho_{11}\vect{v}_{s,h}^{i+1} +\rho_{12}\vect{v}_{f,h}^{i+1}}{\vect{v}_{s,h}^{i+1}} +\Thprod{\rho_{12}\vect{v}_{s,h}^{i+1}+\rho_{22}\vect{v}_{f,h}^{i+1}}{\vect{v}_{f,h}^{i+1}} 
		\\
		+ \Thprod{s_{0}p_h^{i+1}}{p_h^{i+1}} + \frac{\Delta t}{2} \Thprod{\frac{\eta}{\kappa} \vect{v}_{f,h}^{i+1}}{\vect{v}_{f,h}^{i+1}} 
		\\
		+ \frac{\Delta t}{2}  \pThprod{\tau_{s} ({P_{\widehat{V}_s}} \vect{v}_{s,h}^{i+1} - \widehat{\vect{v}}_{s,h}^{i+1})}{P_{\widehat{V}_s}\vect{v}_{s,h}^{i+1}-\widehat{\vect{v}}_{s,h}^{i+1}}
		\\
		+ \frac{\Delta t}{2} \pThprod{\tau_{f} (p_{h}^{i+1} - \widehat{p}_{h}^{i+1})}{p_h^{i+1}-\widehat{p}_h^{i+1}}&=0 . 
	\end{align*}
	%
	Then, using the positivity of $\Ael$ assumption \eqref{eq:A-positive},  and the assumption on the density coefficientes \eqref{eq:rho-coercivity}, we obtain that 
	$\vect{v}_{s,h}^{i+1}= \vect{v}_{f,h}^{i+1} =   0$, $ (\sigma_h^{i+1} + \alpha p_h^{i+1}) = 0, $ and 
	%
	\begin{align*}
		\quad \pThprod{\tau_{s} ({P_{\widehat{V}_s}} \vect{v}_{s,h}^{i+1} - \widehat{\vect{v}}_{s,h}^{i+1})}{P_{\widehat{V}_s}\vect{v}_{s,h}^{i+1}-\widehat{\vect{v}}_{s,h}^{i+1}} = 0, \quad 
		\pThprod{\tau_{f} (p_{h}^{i+1} - \widehat{p}_{h}^{i+1})}{p_h^{i+1}-\widehat{p}_h^{i+1}} = 0. 
	\end{align*}
	%
	Therefore, $\widehat{\vect{v}}_{s,h}^{i+1} = 0$ and $p_h^{i+1} = \widehat{p}_h^{i+1}$ on $\partial \mathcal{T}_h$. Moreover, replacing $\vect{v}_{s,h}^{i+1}= \vect{v}_{f,h}^{i+1} =   0$, and $p_h^{i+1} = \widehat{p}_h^{i+1}$ in the third equation of the system above, and using integration by parts gives $\grad p_h^{i+1} = 0$, so $p_h^{i+1}$ is constant on $\Omega$. Since $p_h^{i+1} = \widehat{p}_h^{i+1}$ on $\partial \mathcal{T}_h$ and $\widehat{p}_h^{i+1}|_{\Gamma_p} = 0$, $p_h^{i+1} = 0$ on $\Omega$. Then, $\sigma_h^{i+1} = 0$ and $\widehat{p}_h^{i+1} =0$. Therefore, \eqref{eq:fully-discrete-cn-eqs} is well-posed.
\end{proof}

\begin{lemma}
	The system \eqref{eq:fully-discrete-cn-eqs} is well-posed.
\end{lemma}

\begin{proof}
	We first note that \eqref{eq:fully-discrete-cn-eqs} is a square linear system; thus, the uniqueness of solutions implies the existence of solutions. 
	Suppose then that $(\sigma_h^i, \vect{v}_{s,h}^i, \widehat{\vect{v}}_{s,h}^i, \vect{v}_{f,h}^i, p_h^i, \widehat{p}_h^i)$ and $g^i$, $g^{i+1}$, $\vect{f}^{i}$, $\vect{f}^{i+1}$ vanish. Then, \eqref{eq:fully-discrete-cn-eqs} reduces to 
	%
	% \begin{subequations}
		\begin{align*}
			% 
			\Thprod{\Ael(\sigma_{h}^{i+1} + \alpha {p}_{h}^{i+1} \Id)}{r} +\frac{\Delta t}{2} (\Thprod{\divergence r}{\vect{v}_{s,h}^{i+1}} -\pThprod{\widehat{\vect{v}}_{s,h}^{i+1}}{r \vect{n}})& = 0 ,
			\\ 
			% 
			\Thprod{\rho_{11}{\vect{v}}_{s,h}^{i+1} +\rho_{12}{\vect{v}}_{f,h}^{i+1}}{\vect{w}_{s}} + \frac{\Delta t}{2}\left(\Thprod{\sigma_{h}^{i+1}}{\grad \vect{w}_{s}} 
			- \pThprod{\widehat{\sigma_{h} \vect{n}}^{i+1} }{\vect{w}_{s}}\right)& = 0, 
			\\
			% 
			\Thprod{\rho_{12}{\vect{v}}_{s,h}^{i+1}+\rho_{22}{\vect{v}}_{f,h}^{i+1}}{\vect{w}_{f}} + \frac{\Delta t}{2}\left(\Thprod{\frac{\eta}{\kappa} \vect{v}_{f,h}^{i+1}}{\vect{w}_{f}} 
			- \Thprod{p_{h}^{i+1} }{\divergence\vect{w}_{f}} +\pThprod{\widehat{p}_{h}^{i+1}}{\vect{w}_{f}\cdot\vect{n}}\right) &= 0, 
			\\
			% 
			\Thprod{s_{0}{p}_{h}^{i+1}}{q} + \Thprod{\Ael({\sigma}_{h}^{i+1}+\alpha {p}_{h}^{i+1}\Id)}{\alpha q\Id} + \frac{\Delta t}{2} \left(-\Thprod{\vect{v}_{f,h}^{i+1}}{\grad q} 
			+ \pThprod{\widehat{\vect{v}_{f,h}\cdot \vect{n}}^{i+1}}{q}\right) &= 0,
			\\
			% 
			\frac{\Delta t}{2}\pThprod{\widehat{\sigma_h \vect{n}}^{i+1} }{\widehat{\vect{w}}_{s}}=
			\frac{\Delta t}{2}\pThprod{\sigma_h^{i+1} \vect{n} - \tau_{s} (P_{\widehat{V}_s} \vect{v}_{s,h}^{i+1} - \widehat{\vect{v}}_{s,h}^{i+1})}{\widehat{\vect{w}}_{s}} &= 0, 
			\\
			% 
			\frac{\Delta t}{2}\pThprod{\widehat{\vect{v}_{f,h}\cdot \vect{n}}^{i+1})}{\widehat{q}}=
			\frac{\Delta t}{2}\pThprod{\vect{v}_{f,h}^{i+1}\cdot \vect{n} + \tau_f (p_h^{i+1} - \widehat{p}_h^{i+1})}{\widehat{q}} &= 0 .
		\end{align*}
		% \end{subequations}
	%
	By taking the test functions $r = \sigma_h^{i+1}$, $\vect{w}_{s} = \vect{v}_{s,h}^{i+1}$, $\vect{w}_{f} = \vect{v}_{f,h}^{i+1}$, $q = p_h^{i+1}$, $\widehat{\vect{w}}_{s} = \widehat{\vect{v}}_{s,h}^{i+1}$, $\widehat{q}=\widehat{p}_{h}^{i+1}$ in the system above and adding all the equations, we obtain
	%
	\begin{align*}
		\Thprod{\Ael({\sigma}_{h}^{i+1} + \alpha {p}_{h}^{i+1} \Id)}{{\sigma}_{h}^{i+1} + \alpha {p}_{h}^{i+1} \Id)}
        \\
		+ \Thprod{\rho_{11}\vect{v}_{s,h}^{i+1} +\rho_{12}\vect{v}_{f,h}^{i+1}}{\vect{v}_{s,h}^{i+1}} +\Thprod{\rho_{12}\vect{v}_{s,h}^{i+1}+\rho_{22}\vect{v}_{f,h}^{i+1}}{\vect{v}_{f,h}^{i+1}} 
		\\
		+ \Thprod{s_{0}p_h^{i+1}}{p_h^{i+1}} + \frac{\Delta t}{2} \Thprod{\frac{\eta}{\kappa} \vect{v}_{f,h}^{i+1}}{\vect{v}_{f,h}^{i+1}} 
		\\
		+ \frac{\Delta t}{2}  \pThprod{\tau_{s} ({P_{\widehat{V}_s}} \vect{v}_{s,h}^{i+1} - \widehat{\vect{v}}_{s,h}^{i+1})}{P_{\widehat{V}_s}\vect{v}_{s,h}^{i+1}-\widehat{\vect{v}}_{s,h}^{i+1}}
		\\
		+ \frac{\Delta t}{2} \pThprod{\tau_{f} (p_{h}^{i+1} - \widehat{p}_{h}^{i+1})}{p_h^{i+1}-\widehat{p}_h^{i+1}}&=0 . 
	\end{align*}
	%
	Then, using the positivity of $\Ael$ assumption \eqref{eq:A-positive},  and the assumption on the density coefficientes \eqref{eq:rho-coercivity}, we obtain that 
	$\vect{v}_{s,h}^{i+1}= \vect{v}_{f,h}^{i+1} =   0$, $ (\sigma_h^{i+1} + \alpha p_h^{i+1}) = 0, $ and 
	%
	\begin{align*}
		\quad \pThprod{\tau_{s} ({P_{\widehat{V}_s}} \vect{v}_{s,h}^{i+1} - \widehat{\vect{v}}_{s,h}^{i+1})}{P_{\widehat{V}_s}\vect{v}_{s,h}^{i+1}-\widehat{\vect{v}}_{s,h}^{i+1}} = 0, \quad 
		\pThprod{\tau_{f} (p_{h}^{i+1} - \widehat{p}_{h}^{i+1})}{p_h^{i+1}-\widehat{p}_h^{i+1}} = 0. 
	\end{align*}
	%
	Therefore, $\widehat{\vect{v}}_{s,h}^{i+1} = 0$ and $p_h^{i+1} = \widehat{p}_h^{i+1}$ on $\partial \mathcal{T}_h$. Moreover, replacing $\vect{v}_{s,h}^{i+1}= \vect{v}_{f,h}^{i+1} =   0$, and $p_h^{i+1} = \widehat{p}_h^{i+1}$ in the third equation of the system above, and using integration by parts gives $\grad p_h^{i+1} = 0$, so $p_h^{i+1}$ is constant on $\Omega$. Since $p_h^{i+1} = \widehat{p}_h^{i+1}$ on $\partial \mathcal{T}_h$ and $\widehat{p}_h^{i+1}|_{\Gamma_p} = 0$, $p_h^{i+1} = 0$ on $\Omega$. Then, $\sigma_h^{i+1} = 0$ and $\widehat{p}_h^{i+1} =0$. Therefore, \eqref{eq:fully-discrete-cn-eqs} is well-posed.
\end{proof}

\end{document}
